\section{Working schemes and uniform comparison of diagrams}
\label{lev:section-working-schemes}

We isolate the conditions used when the preceding constructions are applied
to finite systems of models.  The definition includes the hypotheses that
must be checked for a new system; it does not infer them merely from
categoricity of its vertices.

\subsection{Systems and the working package}

\begin{convention}[Closure below a level]
\label{lev:level-convention}
An eligible level is an infinite cardinal belonging to the class of
cardinals on which the particular construction under consideration has
been established.  Except when explicitly discussing preliminary levels,
we use strong limit levels of countable cofinality.  A list is
\emph{bounded at $\mu$} if its cardinality is strictly less than $\mu$.
If $u$ is a cardinal-valued function, strict closure under $u$ means
\[
  u(\theta)<\mu
  \qquad\text{for every infinite cardinal }\theta<\mu.
\]
Bounds may include fixed parameters, may be increased to infinite
cardinals, and may be replaced by monotone majorants.
\end{convention}

A system has finitely many vertices, each a $\K$-model, and specified
maps between specified vertices.  We code it as a many-sorted structure,
with one sort for each occurrence of a vertex and function symbols for
the specified maps.  An element occurring in two vertices therefore
has two occurrences in the code, related by the specified map when
appropriate.  Distinguished lists are coded by names.  A system
embedding consists of commuting embeddings on its vertices, preserving
these names.  An \emph{arrow} is such an embedding whose vertex maps are
strong, together with any additional condition stated for the scheme.
The geometric application will have one such additional condition.
No condition on intersections of images is implicit in this definition.
An isomorphism means a structural isomorphism onto the entire system.

\begin{definition}[Working package]
\label{lev:working-package}
Fix an unbounded definable class $S$ of eligible levels, a finite system
shape, and a set-sized finite-support term template $F$.  A
\emph{working scheme} consists of classes $\mathcal C_\mu$, for $\mu\in S$,
arrows between their objects, and comparisons $D$ of bounded lists,
with the following properties.  Every vertex of an object in
$\mathcal C_\mu$ has cardinality $\mu$.
\begin{enumerate}
\item \emph{Structural arrows and unions.}
Arrows include all structural isomorphisms, compose, and are invariant
under isomorphisms.  They are coherent: if $X$ is a subsystem of $Y$,
and both inclusions into $Z$ are arrows, then $X\to Y$ is an arrow.
For an increasing nonempty chain of objects at level $\mu$, whose
union has total cardinality $\mu$, the union is an object at that level
and the inclusions into it are arrows.  If all members of the chain
map by the inclusions into one object $Z$ at an eligible level
$\nu\geq\mu$, then the inclusion of the union into $Z$ is an arrow.

\item \emph{Presentations and mixed downward closure.}
For every well-order $I$ of eligible cardinality $\mu$, $F(I)$ is an
object of $\mathcal C_\mu$.  Its elements are represented by labels
applied to finite increasing lists from $I$; equality, the symbols of
the system structure, and the specified maps are evaluated from the
labels and the finite order and equality patterns.  Substitution of
positions along an order embedding induces an arrow, including when
the source and target levels differ.  Every object at a level is
isomorphic to a well-ordered presentation.  If $\nu\leq\mu$ are eligible
and $A$ is a set of at most $\nu$ vertex occurrences in an object of
level $\mu$, there is an arrow inclusion of an object of level $\nu$
containing $A$.  Fixed distinguished data are included in this request.

\item \emph{Diagrams, homogeneity, and saturation.}
For each index set of bounded cardinality there is an equivalence
relation, denoted by equality of $D$-diagrams, on the corresponding
lists in all eligible objects large enough for that index set.
It includes the atomic diagram of the entire coded system, respects
restrictions and reindexing, and sends every list to a $D$-equal list
under every arrow, including mixed-level arrows.  At each level,
two bounded lists have the same diagram if and only if a structural
isomorphism between their objects takes one list to the other.
The number of diagrams at a fixed arity has a set-sized bound
independent of the level.  A possible joint diagram of bounded length
extending the diagram of an actual bounded list is realized over that
list at its level.  Here ``possible'' includes a diagram witnessed at
any larger eligible level, provided the total requested length is
bounded at the receiving level.  On term lists diagrams are natural:
they depend only on labels and the order and equality pattern of all
positions used, including in comparisons between eligible levels.

\item \emph{Exhaustion criterion for arrows.}
Suppose a map is obtained from compatible diagram matches on increasing
bounded lists which exhaust its domain object.  Include in these lists,
for every vertex $V$, nested enumerations of strong submodels
$V_n\leq_\K V$ of sizes below the source level, with union $V$.
The image of each such enumeration must enumerate a strong submodel
of the corresponding target vertex.  Under these conditions the map
is an arrow.  Component strongness follows from the chain union and
common-upper-bound axioms of $\K$; an additional arrow condition, when
present, must be proved for this construction separately.
\end{enumerate}
The empty list is included in the third condition, and has one
diagram.  Consequently all objects at any one eligible level are
isomorphic.
\end{definition}

The mixed downward assertion in the second condition usually needs no
separate construction.  In a chart for the large object, choose finite
supports for the requested occurrences.  Their union has size at most
$\nu$; enlarge it to a suborder of size $\nu$ and use the presentation
arrow.  If an already given small object is also to be included, put
supports for all its occurrences into the request and use coherence.

The assertion about images of small strong submodels in the fourth
condition is also a precise obligation, rather than a claim that
arbitrary structural embeddings are strong.  At one level it follows
when the bounded enumeration is carried by a diagram match, since that
match is implemented by an isomorphism of the containing objects.
For a mixed-level use, first place the smaller containing object in a
large presentation by an order arrow.  Diagram preservation and an
isomorphism at the receiving level then show that the matched small
submodel has strong image.  Finally, the union of these strong images
is strong in the target by the common-upper-bound clause.  This proves
the componentwise part of the criterion in precisely the situations
in which it will be used.

\subsection{Refining the eligible levels}
\label{lev:refining-levels}

The later scheme constructions refine the class of levels on which the working
package is available. The following hypothesis ensures that the required
refinements still leave an unbounded supply of levels.

Here and below every bound being closed under is given by a fixed
set-theoretic definition with set parameters.  To see the elementary
closure construction explicitly, start above a prescribed cardinal and
choose strictly increasing infinite cardinals $\kappa_n$ such that
\[
 \kappa_{n+1}>2^{\kappa_n}
 \quad\text{and}\quad
 \kappa_{n+1}>
 \sup\{u_i(\theta):\theta\leq\kappa_n,\ i\in A\},
\]
where $A$ is the fixed set of functions to be accommodated.  Replacement
makes the displayed supremum a set-sized ordinal.  The cardinal
$\mu=\sup_{n<\omega}\kappa_n$ has countable cofinality, is a strong limit,
and is strictly closed under every $u_i$.

There are two distinct uses of this observation.  If a theorem proves
its conclusion at every such closure point, it establishes eligibility
of $\mu$.  If eligibility is already restricted to an unbounded class
$S$, closing under the least member of $S$ above $\theta$ ensures only
that $S\cap\mu$ is unbounded in $\mu$; it does not by itself prove
$\mu\in S$.  Membership at the limiting cardinal must come from the
relevant theorem or from an already established closure property of
$S$.  In particular, if $S$ is unbounded and closed under suprema of
strictly increasing countable sequences, we may choose the $\kappa_n$
inside $S$ in the construction above.  The resulting $\mu$ then belongs
to $S$ and has cofinally many smaller members of $S$.  This is the
meaning of requiring sufficiently many smaller eligible levels; an
arbitrary unbounded class of cardinals is not silently assumed to have
this property.

\begin{convention}[Supply of eligible levels]
\label{lev:eligible-subclass}
For the generic scheme refinements below, we use the following
\emph{eligible-subclass hypothesis} whenever an unbounded class of
refined levels is asserted.  The input eligible class $S$ must contain
a fixed definable unbounded subclass $T$ closed under suprema of
strictly increasing countable sequences.  In addition to the input
scheme's individual conditions, this is a hypothesis on the supply
of eligible levels; arbitrary unboundedness of $S$ does not imply it.
\end{convention}

This hypothesis persists through every finite list of the stated
refinements.  For a fixed definable bound function, choose successive
members of $T$ beyond all its values on the ordinals below the previous
member, and take their countable supremum.  This gives arbitrarily
large members of $T$ strictly closed under that bound.  The intersection
with the bound-closure class is still closed under strictly increasing
countable suprema.  A threshold only removes an initial segment.
For a limit-point refinement of a current class containing such a $T$,
use $\{\mu\in T:\sup(T\cap\mu)=\mu>0\}$; it is unbounded and has the same countable closure,
and is contained in the required limit-point class.  Repeat for the
finite refinement list, including the finitely many bound functions
then required.  Thus an unbounded subclass with the same property
remains available, including for larger auxiliary-level choices.

\paragraph{Advance notation for common levels.}
\label{lev:closure-preview}
For a definable class $A$ of ordinals, write
\[
 \Lim(A)=\{\mu:\sup(A\cap\mu)=\mu>0\}.
\]
Here is the hierarchy that will organize the two applications of
Section~\ref{close:section-uniform}.  That section constructs one fixed
bound majorant $\Phi$ for each application.  Its initial class $G_0$
consists of the strong limit cardinals $\mu$ of countable cofinality with
$\Phi(\alpha)<\mu$ for every ordinal $\alpha<\mu$; subsequent classes are
\[
 G_{j+1}=G_j\cap\Lim(G_j)\quad(j<\omega),
 \qquad G_\omega=\bigcap_{j<\omega}G_j.
\]
This is advance notation: the definition of $\Phi$ and the proof that
each $G_j$ is unbounded and closed under strictly increasing countable
suprema are given in that section.  A member of $G_{j+1}$ is thus both
a member of $G_j$ and a limit of smaller members of $G_j$.

In those applications, the input class contains $G_j$ for some finite
$j$.  This subclass serves as $T$, and the finite refinements retain a
later $G_k$.  These subclasses supply the required levels; they are not inserted as
additional class parameters into the definitions of bounds. Those
definitions, given in Section~\ref{close:section-uniform}, use fixed
formulas for the required working assertions and for the finite
refinements already described.  Without the eligible-subclass hypothesis, the generic
proofs assert their conclusions at levels satisfying their displayed
bound and limit-point requirements, but do not assert that those
levels are unbounded.

\subsection{Directed unions and rich orders}

\begin{lemma}[Directed unions and padded presentations]
\label{lev:directed-unions}
Assume the structural union and presentation conditions of
Definition~\ref{lev:working-package} at a level $\mu$.
Every nonempty directed family of inclusion subobjects at that level,
whose union has total cardinality $\mu$, has an object as its union,
with arrow inclusions and the common-upper-bound conclusion, including
into an eligible larger level.

Consequently $F(I)$ is an object whenever $I$ has eligible cardinality
$\mu$ and contains a well-ordered subset of cardinality $\mu$.
We call such an order \emph{padded}.  Order embeddings between padded
orders induce arrows, also between different eligible levels.
If the diagram conditions of the package hold, diagrams on bounded
term lists in padded orders are natural as well.
\end{lemma}

\begin{proof}
We first prove the directed assertion by induction on the cardinality
of an indexing set.  A finite directed family has a member containing
all its members.  A countable directed family has a cofinal increasing
sequence, obtained by successively choosing an upper bound for the
previous choice and the next enumerated member.  The chain assertion
gives its union and its common-upper-bound property.

For an uncountable indexing cardinal $\tau$, fix a binary operation on
the indexing set selecting an upper bound for each pair.  Write the
indexing set as an increasing union of sets of size less than $\tau$,
and close each under this operation.  This closure adds only countably
many stages and preserves the bound of being smaller than $\tau$.
The sets can be chosen nested, and their union is the whole indexing
set.  By the inductive assertion each corresponding smaller directed
family has an object as its union.  That union still has cardinality
$\mu$: it contains a member of cardinality $\mu$ and is contained in
the given total union of cardinality $\mu$.  The common-upper-bound
part of the inductive assertion makes these unions a chain of arrows.
The chain assertion now gives the total union.  Applying the same
common-upper-bound clauses to a specified ambient object proves the
last assertion, including the mixed-level version.

Fix a well-ordered subset $J_0\subseteq I$ of size $\mu$.  Let
$\mathcal J$ consist of all well-ordered subsets of $I$ of size $\mu$.
A union of two well-ordered subsets is well-ordered: in any nonempty
subset take the smaller of the minima supplied by its two nonempty
intersections.  Thus $\mathcal J$ is directed under inclusion.  Every
finite subset of $I$ is contained in a member, namely its union with
$J_0$.  The finite-support description of the template identifies
\[
                  F(I)=\bigcup_{J\in\mathcal J}F(J).
\]
The union has cardinality $\mu$, and the directed assertion applies.
For an order embedding $I\to I'$, first restrict to well-ordered
suborders of the source.  Their images are well-ordered; if necessary,
enlarge each image by a padding subset in the target.  The known order
arrows and the common-upper-bound assertion give the induced map on
the total union as an arrow.

Finally, two occurrences of the same bounded support pattern in
orders can be identified by amalgamating the orders over that
suborder.  Such an amalgamation is obtained cut by cut: retain the
common suborder and put the points from each of the two orders in
their prescribed cuts, keeping their own orders within each cut.
Append a well-order of sufficiently large eligible cardinality to
make the amalgam padded.  The two order arrows identify the term
lists there.  Diagram preservation proves their naturalness.
\end{proof}

When the chain assertion for a newly defined scheme is still being
proved, the lemma is used only for an underlying scheme whose chain
assertion has already been established.  Directed unions are not used
to presume the assertion under proof.

\begin{lemma}[Rich padded orders]
\label{lev:rich-orders}
Let $\eta$ be infinite.  There is a dense linear order $Q_\eta$ without
endpoints of cardinality $2^\eta$ such that every cut given by at most
$\eta$ parameters on either side is filled.  Equivalently, every
extension of an order pattern with at most $\eta$ old and at most
$\eta$ added positions can be realized over its old positions.
If $\mu$ is a strong limit cardinal and $\eta<\mu$, there is a padded
order with these properties of cardinality $\mu$.  No restriction on
$\cf(\mu)$ is needed for this assertion.
\end{lemma}

\begin{proof}
Put $\xi=2^\eta$, so that $\xi^\eta=\xi$ and $\eta^+\leq\xi$.
Starting with an order of size $\xi$, repeatedly fill every cut
specified by disjoint sets $A<B$ with $|A|+|B|\leq\eta$, allowing an
empty side.  At each stage there are at most $\xi^\eta=\xi$ requests;
they can be filled in an order extension of size $\xi$.  Perform
$\eta^+$ stages and take unions at limits.  A set of at most $\eta$
parameters in the final union occurs in one earlier stage, since
$\eta^+$ is regular.  Its requested cut is filled at the next stage.
The order is dense and without endpoints because requests with at
most two parameters were among those filled.  To realize an arbitrary
added pattern, well-order its at most $\eta$ new positions and insert
them successively.  Each insertion asks for a cut over at most
$\eta$ old or previously inserted positions.

For the second assertion, use the lexicographic sum
\[
                I_{\mu,\eta}
                =\sum_{\alpha<\mu\cdot\eta^+}Q_\eta,
\]
where the index is ordinal multiplication.  Its cardinality is
$\mu$, its index has cofinality $\eta^+$, and choosing one point in
each summand gives a well-ordered subset of cardinality $\mu$.
Here $|Q_\eta|=2^\eta<\mu$.

Consider a cut $A<B$ with at most $\eta$ parameters.  If $B$ is empty,
the indices of the summands meeting $A$ are bounded in
$\mu\cdot\eta^+$, and a later summand supplies a point above $A$.
If $B$ is nonempty, let $\beta$ be the least index of a summand
meeting $B$.  All points of $A$ lie in summands of index at most
$\beta$.  In the summand at $\beta$, fill the cut between
$A\cap Q_\eta^{(\beta)}$ and $B\cap Q_\eta^{(\beta)}$; the first
set is allowed to be empty.  The resulting point fills the original
cut.  The same successive-insertion argument realizes added patterns.
\end{proof}

\subsection{The initial package and named data}

\begin{proposition}[The package for single models]
\label{lev:single-model-package}
At the sufficiently closed levels supplied by
Proposition~\ref{lev:closure-categoricity}, the single $\K$-models,
with strong maps as arrows and the original term presentation $E$,
form a working scheme.  The levels can be refined to those required
in Theorem~\ref{lev:amalgamation} without losing this conclusion.
\end{proposition}

\begin{proof}
The structural conditions are the AEC axioms, including their
common-upper-bound clauses.  The presentation and its mixed order
arrows are supplied by Theorem~\ref{found:order-template}.
Categoricity at the specified levels identifies every model with
each presentation of the same size; supports in such a presentation
give mixed downward closure as described after
Definition~\ref{lev:working-package}.

Proposition~\ref{lev:stabilization} gives the bounded diagrams, their
uniform cardinal bounds, restriction compatibility, and preservation
under strong maps.  Proposition~\ref{lev:closure-categoricity} gives
isomorphism at equal bounded diagrams.  For saturation, let $a$ be
an actual tuple and let $(b,c)$ witness a possible joint diagram,
with $D(a)=D(b)$.  Choose an infinite $\sigma<\mu$ bounding both
lists.  The closure requirements provide a categoricity cardinal
$\kappa<\mu$ above $\sigma$ and the stabilization thresholds for
these arities.  Put $a$ in a strong $\kappa$-submodel of its actual
model.  Put $(b,c)$ in a strong $\kappa$-submodel of its witness,
first embedding the witness into a larger eligible presentation if
needed.  Stabilization identifies the $\kappa$-comparison classes
of $a$ and $b$, so a strong map from the second submodel to the
first takes $b$ to $a$.  Its image of $c$ realizes the required
joint diagram.  Naturality follows from the order
arrows and order amalgamation.  Images of small strong submodels are
strong by the isomorphism argument following
Definition~\ref{lev:working-package}; taking their increasing unions
proves the exhaustion criterion.  Finally, imposing the additional
closure bounds of Theorem~\ref{lev:amalgamation} retains each of
these already established assertions.  Those additional bounds
provide the amalgamation and comparison facts needed later; they
are not extra axioms of a working scheme.
\end{proof}

\begin{proposition}[Naming bounded data]
\label{lev:named-package}
Let $\mathcal C$ be a working scheme, and fix a bounded list $c$ and
its diagram, witnessed at some eligible level.  After increasing the
fixed lower threshold, the objects with a named realization of this
diagram, and arrows preserving the names, form a working scheme.
This includes naming enumerations of fixed small strong models in
the vertices, with their specified maps.  The additional lower bound
can be chosen in terms of the original scheme and the cardinality
of the named list, independently of the level at which that list
was first witnessed.
\end{proposition}

\begin{proof}
Let $\beta$ be an infinite bound for the length of $c$, and choose
one eligible $\alpha>|F|+\beta$.  The diagram of $c$ has a witness
at $\alpha$.  If its original witness is at a larger level, use
saturation at $\alpha$ with the empty tuple as base.  If its witness
is at a smaller level, first extend that object to a presentation
at a level at least $\alpha$ using an order arrow.  This gives a
choice of $\alpha$ uniform over all diagrams of the specified arity.
Choose a well-ordered chart for this new witness and fix
representations of $c$.  Let $J$ be the union of their finite supports, with its inherited
well-order.  In the presentations $F(J+H)$ name the terms representing
$c$ on $J$.  Regard this as a template in the variable order $H$:
any term using positions from both parts becomes a label recording
its original label and its finitely many fixed positions from $J$,
applied to its remaining positions from $H$.  Terms supported wholly
on $J$ are labels of arity zero.  The new collection of labels has
cardinality at most $|F|+|J|+\aleph_0$ and is fixed independently of
the level.  Order maps on $H$, extended by the identity on $J$,
induce arrows preserving the names.

Naturality gives the prescribed diagram on the named list in every
sufficiently large such presentation.  Homogeneity for $c$ shows
that every object with such a named realization is isomorphic to one
of these presentations.
Define its diagram on a further list $a$ to be the old diagram on
the combined list $(c,a)$.  The old package gives exactly restriction
compatibility, the uniform bound, preservation, homogeneity, and
saturation for these diagrams.  The same support construction gives
mixed downward closure, now including the named terms.  Chain unions
and the exhaustion criterion retain the names, so they are inherited.

For the final assertion suppose a sublist enumerates
$M\leq_\K V$ in the original witness, which may be at a level
larger than $\alpha$.  Embed that original system and the new
level-$\alpha$ witness into objects of a common larger eligible
level.  Mixed preservation and homogeneity there give an
isomorphism matching their named lists.  The original copy of $M$
has strong image in that large vertex.  This image is exactly the
substructure enumerated inside the embedded level-$\alpha$ vertex,
so coherence makes it strong in that vertex as well.  It is a
$\K$-model because its complete atomic enumeration is isomorphic
to $M$.  Thus strong naming already holds in the uniformly small
witness.  Above $\alpha$, order arrows and the same homogeneity
argument preserve it.  The argument applies to every named
vertex model and specified map.
Thus the new presentations satisfy the intended strong naming
condition, not just the atomic diagram of the enumeration.
\end{proof}

\subsection{Two routes to working diagrams}

The next statement records the exact diagram input that may be used
for a new scheme.  In its first route the role of
Theorem~\ref{found:reversal} is an explicit hypothesis; it is not
assumed for a new category without proof.

\begin{proposition}[Obtaining the diagram clause]
\label{lev:scheme-diagram-routes}
Suppose a finite-system scheme has the structural and presentation
conditions of Definition~\ref{lev:working-package}, and is categorical
at its eligible levels.  For an unbounded supply after refinement,
assume the eligible-subclass hypothesis of
Convention~\ref{lev:eligible-subclass}.  Either of the following inputs
supplies the diagram clause, after the indicated closure refinements.
\begin{enumerate}
\item For each infinite $\sigma$, at all sufficiently large eligible
levels, the relation of being sent to another list by an arrow is
reversible on lists of cardinality at most $\sigma$, and their orbit
sets have a bound depending only on $\sigma$ and the fixed scheme.
In this case take stabilization bounds and require enough smaller
levels below the receiving level for the bounded arities in use.

\item Same-level arrows preserve isomorphism orbits of bounded lists,
and comparisons of such orbits in well-ordered charts are independent
of the eligible level.  More precisely, for each arity, equivalence
of two term patterns by an isomorphism in their respective charts is
the same relation on patterns whenever both patterns are bounded at
the levels under comparison.
\end{enumerate}
The componentwise exhaustion condition follows whenever arrows have
only componentwise strongness as their additional structural content.
For any other arrow condition its exhaustion assertion remains a
separate obligation.
\end{proposition}

\begin{proof}
For the first route, fix an arity bound $\sigma$.  At a sufficiently
large eligible cardinal $\kappa$, arrow comparison on lists of this
arity is reflexive and transitive, and reversal makes it symmetric.
A list in a larger object can be put into a level-$\kappa$ subobject.
Two choices of that subobject are included in a third by mixed
downward closure and coherence.  Their arrow classes therefore agree.
This defines the level-$\kappa$ comparison in all larger objects and
shows that it is preserved under mixed arrows.

If $\gamma\geq\kappa$ is another sufficiently large eligible cardinal,
there is consequently a projection from the $\gamma$ comparison
classes to the $\kappa$ comparison classes.  This projection is onto:
an object at level $\kappa$ embeds into one at level $\gamma$ by
identifying it with a chart and using an order arrow.  The projections
compose.  Let $b$ bound the numbers of classes.  If strict refinement
continued arbitrarily high, recursively choose $b^+$ many successive
strict refinements and then an eligible cardinal above all the
chosen levels.  At each strict refinement some pair of classes at
this final level has equal images before the refinement and unequal
images after it.  A fixed pair can witness this only once, since
unequal images cannot merge in a refinement.  There are at most
$b^2=b$ pairs, a contradiction.  Thus the comparisons stabilize.

Their stabilized equality is canonical, and comparisons at different
arities commute with restriction: raise the cardinal beyond the
stabilization thresholds for both arities and restrict an arrow
implementing the longer comparison.  Atomic diagrams are preserved
by arrows.  A prescribed bounded extension is realized over an actual
list in a sufficiently high same-level object by placing a witness
and the actual list into objects of one common eligible size and
using the reversible comparison to map the witness list to the actual
one.  This gives the joint extension as well.

At a receiving level $\mu$ of countable cofinality, close below $\mu$
under these thresholds and under choices of a sufficiently large
smaller eligible cardinal for each arity.  Enumerate each of two
objects by countably many nested bounded lists.  To extend a match,
choose a smaller eligible cardinal containing all the old and newly
requested entries and above their stabilization thresholds.  Put the
data into subobjects at that cardinal and implement the equal
stabilized diagrams by an arrow in the desired direction.  Alternate
the direction of extension and include the next enumeration list on
each side.  The union is a bijection preserving and reflecting every
atomic formula, hence an isomorphism of systems.  Starting with a
given equal-diagram tuple makes it part of that isomorphism.  The
same smaller-level construction realizes a possible bounded extension
diagram over an actual tuple.  Order arrows and order amalgamation
give naturality.  These arguments use both eligibility of $\mu$ and
the availability of smaller eligible cardinals, as distinguished in
Section~\ref{lev:refining-levels}.

For the second route, instances of the same well-ordered pattern
have the same orbit.  Indeed two well-orders amalgamate over a common
suborder into a well-order: in each cut of the common well-order
concatenate the two remaining well-ordered pieces, then retain the
common points between successive cuts.  Enlarge to an eligible size
if necessary.  The induced order arrows and same-level orbit
preservation identify the two instances.  The assumed independence
of orbit comparison from the level therefore defines diagrams on
well-ordered patterns.  The set of such patterns at a fixed arity
gives a uniform bound.

For a mixed arrow from level $\nu$ to level $\mu$, include the image
of its whole source in a level-$\nu$ subpresentation of the target.
Coherence factors the arrow through that subpresentation.  Same-level
orbit preservation computes the comparison there, and the assumed
agreement of chart comparisons computes the same diagram in the
target.  This proves mixed preservation.  For saturation, take a
well-ordered support pattern witnessing the desired joint diagram
in a larger object.  Its support has cardinality below the receiving
level, so it embeds into a well-order of that level.  The old subtuple
has the diagram of the prescribed actual tuple; an isomorphism at
the receiving level carries the new realization over that tuple.
The orbit hypothesis itself supplies homogeneity.  Naturality on
padded orders follows from Lemma~\ref{lev:directed-unions}.

In either route, diagrams on bounded enumerations of small vertex
submodels have strong images by homogeneity, with the preliminary
order embedding in the mixed-level case.  The AEC union axioms give
the asserted componentwise exhaustion condition.
\end{proof}

\subsection{A set-theoretic code for diagram equality}

The next construction supplies one definition of comparison that can
be used uniformly with a finite descriptor for the scheme.  It does
not define truth for a varying family of proper-class predicates.

A code for a finite-support template consists of its set-sized
finitary language, sorts and specified-map symbols, its labels and
their finite arities, and tables evaluating the function symbols,
equality, and relations on finite ordered support patterns.  A
function table supplies a label and sublist for its output within
the union of its inputs' supports; its arity-zero case is included.
Adjoin labels of arity zero for closed terms if necessary.  This
does not increase the infinite bound on the language and label set.
Fixed names, including
the labels of arity zero in Proposition~\ref{lev:named-package}, are
part of this code.  The usual consistency requirements say that the
evaluated equality is an equivalence relation and a congruence, that
the function tables are total and correctly sorted, and that the
evaluations agree under passage to larger finite support
patterns.  These are tests on sets.  A coherent code constructs a
structure $F(I)$ on any set-sized order $I$ by taking the quotient
of the set of labeled finite lists by the evaluated equality.  This
construction makes sense whether or not $F(I)$ is an object of a
working scheme at that size.

For an arity bound $\sigma$, a \emph{pattern} records an indexed list
of labels, finite increasing support lists for them, and the order
on the union of these supports.  Use canonical support universes
of cardinality at most $\sigma$ and index sets of cardinality at most
$\sigma$.  These codes form a set, denoted by $\mathcal P_F(\sigma)$.
Two codes representing the same terms need not be identified in
advance.  The comparison below handles such duplicate representations.
When names are represented by reserved support positions instead of
labels of arity zero, include their fixed support in every pattern
and every support map; the test bounds must then also dominate its
cardinality.

\begin{proposition}[Uniform coding of working diagrams]
\label{lev:diagram-coding}
There is a fixed first-order set-theoretic definition $D_F$ of an
equivalence relation on term patterns, uniform in a coherent
finite-support template code $F$.  There is a definable cardinal
bound $b_F(\sigma)$ such that, on $\mathcal P_F(\sigma)$, a single
back-and-forth test with bound at least $b_F(\sigma)$ computes $D_F$.
These definitions do not require the system axioms at the cardinals
used for the tests.

Suppose, in addition, that $F$ presents a working scheme at unboundedly
many eligible strong limit levels of countable cofinality.  At every
eligible level $\mu>|F|$ strictly closed under $b_F$, the working
diagram of a bounded term list is exactly its $D_F$-comparison class.
The same definition therefore computes the diagrams of every such
scheme, including the schemes with fixed named data.
\end{proposition}

\begin{proof}
Fix an infinite $\eta$ at least the cardinality of the template code
and of the patterns to be tested.  For two set structures $A,B$ in
its language, form the set of all pairs of tuples of the same index
set of cardinality at most $\eta$.  Index sets can be taken to be
ordinals below $\eta^+$, reindexing concatenations when necessary.
Let $R_0$ consist of the pairs with the same atomic diagram.  Having
defined $R_\alpha$, retain in $R_{\alpha+1}$ just those pairs
$(a,b)\in R_\alpha$ such that every added tuple of cardinality at
most $\eta$ on either side has a matching added tuple on the other
side with the extended pair in $R_\alpha$.  At a limit take the
intersection of all preceding relations.

This decreasing recursion stabilizes.  If its state set has
cardinality $\tau$, there cannot be a strict decrease at every
successor stage below $\tau^+$, since each strict decrease removes
a previously unremoved state.  Equality at one successor stage
gives a fixed point.  Write $R_\infty$ for the stabilized relation.
In particular, from a pair in $R_\infty$ every allowed extension has
a matching extension still in $R_\infty$; quantification over all
ranks is not being exchanged with an existential choice without
this fixed-point argument.  Reflexivity and symmetry follow from
the construction.  Transitivity follows by composing matches:
the composition of two relations having the atomic and extension
properties again has those properties, and is contained in the
largest such relation $R_\infty$.  This also proves invariance under
isomorphisms.

Use orders satisfying Lemma~\ref{lev:rich-orders}, and realize the
two indicated patterns in their term structures.  Declare
$p\equiv_{F,\eta}q$ when their tuple pair belongs to $R_\infty$.
Every support pattern of size at most $\eta$ embeds into such an
order, by the successive insertion assertion in that lemma.

The definition is independent of the rich orders and of the
realizations of an individual pattern.  To prove this, relate two
term tuples when their representations are carried to each other
by an order isomorphism between support sets of size at most $\eta$.
This relation has atomic matching by the definition of the template.
An added tuple has at most $\eta$ new support positions.  Order
richness extends the support map to these positions on the other
side, in either direction.  Thus this relation has the extension
property and is contained in $R_\infty$.  Replacing an instance of
either pattern by another instance and using transitivity proves
independence for the comparison of two possibly different patterns.

If $\eta'\geq\eta$, orders rich for $\eta'$ are also rich for $\eta$,
and their back-and-forth test with the larger moves implies the test
with the smaller moves.  Independence of the chosen orders gives
\[
              \equiv_{F,\eta'}\ \subseteq\ \equiv_{F,\eta}
\]
on patterns in the common domain.  Define $D_F$ to be agreement at
every sufficiently large test bound $\eta$.

Here is both uniform definability and the promised single bound.
In defining the test at $\eta$ it is enough to use rich orders of
cardinality $2^\eta$ on a fixed standard universe, whose existence
was proved in Lemma~\ref{lev:rich-orders}.  All their order relations,
the realizations of the patterns, the associated term structures,
and the states of the recursion are sets.  The recursion up to the
successor cardinal of its state set defines $R_\infty$ by an ordinary
first-order set-theoretic formula.  Independence means that one may
quantify over the rich orders and realizations in that set, without
making a uniform class choice of an order.

For each pair in $\mathcal P_F(\sigma)^2$ which fails some test,
take its least failing infinite cardinal above the fixed input bounds.
This is a uniquely defined ordinal when it exists.  Replacement
collects these witnesses into a set.  Choose $b_F(\sigma)$ above
their supremum, above $\sigma$, and above the template and naming
bounds.  Monotonicity proves
\[
 p\mathrel{D_F}q
 \quad\Longleftrightarrow\quad
 p\equiv_{F,b_F(\sigma)}q,
 \qquad p,q\in\mathcal P_F(\sigma).
\]
The same holds at every larger test bound.  Taking suprema over
smaller arities makes $b_F$ monotone if desired.  Notice that the
cardinality of the set of patterns justifies taking the set of
separation witnesses; it does not assert an a priori numerical
upper bound on their ranks.  A code failing the coherence tests may,
if a total definition on all set codes is desired, be assigned the
discrete comparison and a fixed default bound.  This convention has
no mathematical role for a coherent template.

We now compare with the working diagrams.  First suppose two bounded
patterns have equal working diagrams.  Fix any test bound $\eta$.
Choose an eligible level large enough for the patterns and $\eta$,
and use a rich padded order of that level from
Lemma~\ref{lev:rich-orders}.  Lemma~\ref{lev:directed-unions} makes
its term structure an object of the scheme.  Naturalness transfers
the equal working diagrams to realizations of the two patterns
there.  The package gives an isomorphism taking one realization to
the other.  It passes the back-and-forth test at $\eta$.  Independence
of rich orders now gives $p\equiv_{F,\eta}q$.  Since $\eta$ was
arbitrary, equal working diagrams imply $D_F$-equivalence.

Conversely fix an eligible $\mu$ as in the statement and two objects
of that level with bounded lists whose patterns are $D_F$-equivalent.
Choose countable nested bounded enumerations exhausting both objects.
We construct compatible matches containing these enumerations,
maintaining $D_F$-equivalence of their term patterns.  Suppose the
current lists have been matched, and a next bounded list on one side
is requested.  Let an infinite $\rho<\mu$ bound the cardinality of the combined
old and requested lists.  Choose
\[
          b_F(\rho)\leq\eta<\mu.
\]
This is possible by strict closure.  At the same level $\mu$, form
rich padded presentations for this test.  Realize the current
patterns there.  Naturality and homogeneity for the working diagrams
give isomorphisms from the current objects to these presentations
taking their current lists to the chosen realizations.

Transport the requested list through the first isomorphism.  The
current pair passes the test at $\eta$, so the fixed-point extension
property supplies a matching list in the other rich presentation.
Choose finite-support term representations for the extended lists.
Their patterns pass the $\eta$-test; because
$\eta\geq b_F(\rho)$, they are $D_F$-equivalent.  Transport the
matching list back to the actual target object.  The isomorphisms
used in this transport preserve working diagrams, which were
already shown to imply $D_F$-equivalence.  Transitivity therefore
gives the required $D_F$-equivalence for the extended actual lists,
independently of their chosen representations.

Repeat with the sides reversed and with the next lists in the two
enumerations.  All stages have cardinality below $\mu$; only their
countable union exhausts the objects.  Atomic matching ensures that
the compatible assignments define a well-defined injective map, and
the alternating exhaustion makes it onto every sort.  Every atomic
instance involves finitely many entries and occurs at one stage,
so the resulting bijection is a structural isomorphism.  It carries
the original list to the other list.  The working package therefore
identifies their working diagrams, as required.
\end{proof}

\begin{remark}[What the code does and does not establish]
\label{lev:coding-interface}
The code $F$, its finite shape, and its fixed naming data are set
parameters.  The recursion in Proposition~\ref{lev:diagram-coding}
is a recursion on a set of states of set structures.  Its definition
is the same formula for every such parameter, including templates
arising after a finite number of later constructions.  It therefore
introduces no satisfaction predicate for arbitrary class definitions.
The bound $b_F(\sigma)$ depends on the full fixed descriptor; names
or labels of unbounded size cannot be omitted from that input.

The identification of $D_F$-equivalence with equality of working diagrams
uses an already established working package, including its actual
isomorphism assertion.  It supplies a
uniform way to express that package in subsequent closure tests; it
is not a substitute for proving homogeneity, chain closure, or an
additional arrow condition for a new scheme.  The later simultaneous
closure argument must still prove that its descriptors and its bound
recipes are uniformly set-theoretically definable and that their
requested bounds exist on relevant inputs.
\end{remark}
